\subsubsection{Hierarchy Structure}
\label{subsubsec:system-architecture-hierarchy}

Above the cluster layer, the system deploys a configurable hierarchy whose scopes (cluster, state, nation, etc.) are declared through structured configuration. Each scope specifies a name, identifier, timer interval, wait budget, merge policy, and fan-out quota. Nodes ingest a hierarchy map listing every roster and parent-child relationship so they can derive memberships at runtime. During startup, each node instantiates a per-scope runtime that tracks round queues, round-robin aggregator pools, and the fan-out members appointed to relay models upward~\cite{BCFramework,TowardsFederatedLearningScale,Mitra2023TAHFL}.

High-level aggregation is fully time-driven: every scope owns a wall-clock interval, and when its timer fires the members enqueue a round while briefly pausing local optimization for the configured wait window. The elected aggregator receives the latest artifacts from child scopes via bridge fan-out, deduplicates them, fetches the referenced tensors from IPFS, and merges them according to the scope policy. Only one candidate commits the result per round: the merged tensor is published to IPFS, anchored on the ledger, and the commit event is logged for observability~\cite{BCFramework,TrustDFL}. Peers merely watch the ledger to detect completion and mark the round satisfied~\cite{BCFramework}.

Dissemination back to lower layers follows an asynchronous pull model. Once the wait window elapses---or when a fresh anchor appears early---nodes query the ledger, fetch the referenced tensor, verify its hash, and apply the scope policy (replace, interpolate with a configured $\alpha$, etc.) before resuming local updates. Each scope records the last applied CID so redundant anchors are ignored, allowing nodes to adopt the most recent confirmed checkpoint even while the next one is under construction. Because pauses are brief and higher-level checkpoints are never pushed synchronously, the hierarchy does not block cluster-level SAP yet ensures all members eventually incorporate each scope’s consensus model~\cite{TowardsFederatedLearningScale,BCFramework}.

All experiments instantiate two scopes above the cluster layer: a state level that aggregates clusters within an administrative region and a nation level that aggregates state checkpoints, mirroring typical cross-silo deployments~\cite{Kairouz2021,Mitra2023TAHFL}. The hierarchy remains extensible: operators can declare additional levels whenever topology, policy, or regulatory constraints demand them, letting the same control plane span metropolitan overlays or multi-national consortia~\cite{BCFramework}.